\documentclass{article}

\usepackage[english]{babel}
\usepackage[letterpaper,top=2cm,bottom=2cm,left=2cm,right=2cm]{geometry}
\usepackage{amsmath}
\usepackage{graphicx}
\usepackage[colorlinks=true, allcolors=blue]{hyperref}
\setlength{\parindent}{0pt}

\title{Milestone 1}
\author{Eunice Chen, Muhaimen Shamsi, Dariia Kucheruk, Oleksii Nakhod}
\date{CSC2541H / PCL3107H / PCL3108H, Winter 2026}

\begin{document}
\maketitle

\section*{Part I: Introduction and Background (1 Page Maximum)}

\subsection*{1. Research Question \& Significance}
\textbf{Research question.} Do AD-implicated genes show stronger and more specific signal than non-AD controls in (i) DrugCLIP embedding neighborhoods and (ii) independent PPI topology, and can this be used to prioritize plausible drug candidates?

\textbf{Why this matters for AI in drug discovery.} Drug discovery embedding models are usually evaluated on benchmark task performance, but not always on disease-specific biological validity. Our goal is to test whether embedding-derived structure is aligned with AD biology rather than generic signal. If yes, this supports using embedding + network evidence for disease-focused target prioritization.

\subsection*{2. Context \& Datasets}
\textbf{Context.} We combine representation-learning diagnostics and network medicine validation. We use matched controls and degree-matched nulls to avoid inflated claims from graph hubs.

\textbf{Datasets used.}
\begin{itemize}
    \item DrugCLIP protein embeddings and names (THU-ATOM DrugCLIP data) \cite{adGwas2025, drugclipData2024}.
    \item AD gene list from the Askenazi 2023 supplement (processed into \texttt{data/processed/ad\_genes.csv}).
    \item Allen Human Brain RNA-seq atlas for expression-based AD/control label construction \cite{allenAtlas2012}.
    \item STRING human PPI (v12) for topology and degree-matched controls \cite{string2018, string2019}.
    \item DrugCLIP target ChEMBL human assay tables (from \texttt{targets.zip}) for candidate-drug retrieval.
\end{itemize}

\textbf{Current status before Milestone 2.} We already observe moderate signal: weak embedding-only classification lift, but stronger and statistically significant PPI concentration under degree-matched permutation controls.

\section*{Part II: Proposed Methods and Setup (2 Pages)}

\subsection*{3. Proposed Methods}
\textbf{Method A: AD-vs-control prediction head on embeddings.} We aggregate multiple embedding rows to one gene vector and train a small FFN (768$\rightarrow$64$\rightarrow$1). We compare three ablations: embedding, random-embedding, and label-shuffle.

\textbf{Method B: Embedding-to-network biological validation.} We test whether AD labels are topologically concentrated in PPI with degree-matched permutation nulls. Metrics include AD-AD edge enrichment, AD largest connected component, and AD-neighbor fraction.

\textbf{Method C: Cosine-neighbor/PPI coupling by hop depth.} For $k\in\{1,2,5\}$ cosine neighbors, we compare AD-seed vs control-seed neighbor properties (AD fraction and PPI coupling), then evaluate significance with degree-matched permutation.

\textbf{Method D: Candidate drug retrieval and characterization.} We retrieve potent compounds from DrugCLIP ChEMBL target tables for AD-overlap targets (with alias mapping), rank compounds by AD-target coverage and potency support, then characterize phase/potency/target distributions.

\textbf{Innovation claim.} The novelty is not a new foundation model; it is a disease-specific validation stack that links embedding geometry to independent biological topology using strict controls (label-shuffle, random-embedding, degree-matched nulls).

\subsection*{4. Data and Preprocessing}
\textbf{ID harmonization.} Gene symbols are standardized to uppercase and mapped to UniProt accessions (HGNC mapping) for embedding joins.

\textbf{Leakage prevention.}
\begin{itemize}
    \item Aggregation to one embedding per gene is done \emph{before} train/test splitting.
    \item Splits are fixed by saved split files for fair ablation comparison.
    \item PPI analyses use degree-matched permutation controls to reduce hub-driven artifacts.
\end{itemize}

\textbf{Filtering and thresholds.}
\begin{itemize}
    \item PPI analyses use STRING combined score threshold (default 700 for high confidence).
    \item Drug retrieval keeps potency-supported rows (default: pChEMBL $\geq 6$ or standard value $\leq$ 1000 nM where applicable).
\end{itemize}

\subsection*{5. Information Flow Diagram}
\begin{figure}[h]
    \centering
    \includegraphics[width=\linewidth]{images/information_flow.png}
    \caption{Architecture and information flow used in Milestone 1 experiments.}
\end{figure}

\subsection*{6. Experimental Setup \& Baselines}
\textbf{Primary metrics.}
\begin{itemize}
    \item Predictor: test AUROC, AUPRC, accuracy (multi-seed mean/stdev).
    \item PPI topology: AD-AD enrichment, AD LCC size, AD-neighbor fraction with permutation p-values.
    \item Hop analysis: AD-control delta at $k=1,2,5$ with degree-matched permutation p-values.
\end{itemize}

\textbf{Baselines locked in.}
\begin{itemize}
    \item Random embedding baseline.
    \item Label-shuffle baseline.
    \item Degree-matched permutation null for PPI and hop analyses.
\end{itemize}

\textbf{Current quantitative snapshot.}
\begin{itemize}
    \item Multi-seed embedding AUROC: 0.524 (embedding) vs 0.484 (random embedding) vs 0.516 (label-shuffle).
    \item Degree-matched PPI: AD-AD enrichment $\approx$ 1.29x, p $\approx$ 0.0033; AD-LCC p $\approx$ 0.0033.
    \item Cosine/PPI hops: AD-neighbor fraction is consistently higher for AD seeds; significant under degree-matched permutation in current runs.
    \item Drug retrieval coverage increased from 2 to 5 AD-overlap targets after alias/composite mapping.
\end{itemize}

\subsection*{7. Timeline to Milestone 2}
\textbf{Week of March 8--14, 2026}
\begin{itemize}
    \item Freeze data snapshots and rerun all key analyses at final permutation depth.
    \item Owner: Muhaimen (pipeline + experiment reproducibility), Dariia (run audits).
\end{itemize}

\textbf{Week of March 15--21, 2026}
\begin{itemize}
    \item Expand AD target coverage and finalize candidate-drug characterization tables.
    \item Owner: Eunice (biological curation), Oleksii (data joins and QC).
\end{itemize}

\textbf{Week of March 22--28, 2026}
\begin{itemize}
    \item Final ablation tables/figures and robustness appendix (threshold sweeps, split sensitivity).
    \item Owner: Muhaimen + Dariia.
\end{itemize}

\textbf{Week of March 29--April 4, 2026}
\begin{itemize}
    \item Draft final milestone narrative integrating methods, controls, and biological interpretation.
    \item Owner: Eunice + Oleksii.
\end{itemize}

\textbf{Anticipated bottlenecks.}
\begin{itemize}
    \item Long permutation jobs on large PPI graphs.
    \item Runtime variability from large table parsing and multi-seed reruns.
    \item Limited AD-target overlap in current DrugCLIP target archive.
\end{itemize}

\bibliographystyle{unsrturl}
\bibliography{../loi/loi}

\end{document}
